\section{把公式落实到现有 Python 项目}
\label{code:chapter}

这一节按实际计算顺序阅读现有代码。先理解一个数组代表什么，再看如何更新状态，最后看误差怎样汇总成报告。这里的短代码取自原程序或本讲义附带的教学脚本；为便于阅读而省略上下文的片段不应单独当作完整程序运行。原项目解决大规模实验与报告生成，新增的 \texttt{teaching\_demo.py} 则保留学习所需的核心计算，只依赖 NumPy，默认打印结果，不写文件。

\subsection{文件之间怎样配合}

\begin{center}
\small
\begin{tabularx}{\textwidth}{@{}lY@{}}
\toprule
文件 & 职责与阅读重点\\
\midrule
\texttt{gbm.py} & 生成共同布朗增量；计算精确 GBM、价格 EM、Milstein；独立预试验拟合控制变量；保存收敛统计、终值样本和图。\\
\texttt{nonaffine.py} & 生成相关布朗增量；对 $(X,Y)$ 做 log-EM；比较粗网格与三层参考网格；记录状态与参考敏感性。\\
\texttt{verify.py} & 用确定性输入和另一种写法核对算法，再由保存的终值数组重新算误差，检查统计是否一致。\\
\texttt{build\_report\_data.py} & 从已有 JSON/CSV 生成报告使用的数值宏、八张表和四幅图，并更新 \texttt{numerical\_summary.txt}。\\
\texttt{run\_all.py} & 依次运行上面四个脚本；某一步失败就停止；它本身不编译 \LaTeX。\\
\texttt{README.md} & 记录依赖、实验规模、命令、误差定义和解释边界。\\
\bottomrule
\end{tabularx}
\end{center}

\texttt{results/} 保存计算结果，\texttt{figures/} 保存图，原报告的 \texttt{.tex} 把这些结果排版。因而“重画已有结果”“重新模拟所有路径”“重新编译 PDF”是三个不同操作。原项目的这几个生成脚本使用固定输出文件名；学习时先运行新增教学脚本，可以保留现有实验结果。

\subsection{NumPy 数组：先问每一个轴代表什么}

NumPy 的 \texttt{array} 是规则排列的数。\texttt{shape} 给出每个轴上的长度，轴编号从 0 开始。设本批有 $B$ 条独立路径，最细网格有 $N_f$ 步。GBM 使用“路径在前”的数组：
\begin{lstlisting}[language=Python]
fine = rng.normal(0, math.sqrt(T / nfine), (size, nfine))
w = fine.sum(axis=1)
\end{lstlisting}
这里 \texttt{size} 就是 $B$。\texttt{fine.shape == (B, Nf)}，\texttt{fine[j,k]} 是第 $j$ 条路径、第 $k$ 个小时间段的布朗增量。每一行是一条路径。正态生成器的第二个参数是\emph{标准差}，所以传入 $\sqrt{T/N_f}$，不是方差 $T/N_f$。沿 \texttt{axis=1} 求和，就是把每一行的所有时间增量相加，得到 $W_T$；结果 \texttt{w.shape == (B,)}。

若粗网格有 $N$ 步，且 $r=N_f/N$ 是整数，原代码写作
\begin{lstlisting}[language=Python]
dw = fine.reshape(size, int(n), nfine // int(n)).sum(axis=2)
\end{lstlisting}
\texttt{reshape} 不重新抽随机数，它把原有顺序整理为 $(B,N,r)$。最后一个轴含有同一粗区间内连续的 $r$ 个细增量；沿该轴求和，输出 $(B,N)$。例如单条路径八个细增量为
\[
(a_1,a_2,a_3,a_4,a_5,a_6,a_7,a_8),
\]
取 $N=2$ 时，先分为 $(a_1,\ldots,a_4)$ 与 $(a_5,\ldots,a_8)$，再变为 $(a_1+\cdots+a_4,a_5+\cdots+a_8)$。这是“粗细路径共用同一布朗运动”的代码实现。若另抽一套粗增量，即使分布正确，逐路径的差也不再衡量同一随机驱动下的离散误差。

非仿射程序则使用“时间在前”的数组。两种布局都正确，但求和轴不同：
\begin{center}
\small
\begin{tabularx}{\textwidth}{@{}lYY@{}}
\toprule
操作 & GBM & 非仿射模型\\
\midrule
最细增量形状 & $(B,N_f)$ & $(N_f,B,2)$\\
一个元素 & 第 $j$ 条路径、第 $k$ 步 & 第 $k$ 步、第 $j$ 条路径、第 $q$ 个噪声\\
粗化前重新分组 & $(B,N,r)$ & $(N,r,B,2)$\\
应求和的轴 & \texttt{axis=2} & \texttt{axis=1}\\
粗化后形状 & $(B,N)$ & $(N,B,2)$\\
\bottomrule
\end{tabularx}
\end{center}
不要只背 \texttt{axis=1} 或 \texttt{axis=2}；应先写出形状，再指认“要合并的细时间轴”。类似地，\texttt{sum(axis=0)} 在 $(B,2)$ 上是跨路径汇总，结果中保留两种方法。

\subsection{GBM：从递推式到 \texttt{prod}}

因为 GBM 的系数都正比于价格，价格 EM 和 Milstein 都能写为“旧价格乘一个因子”。原 \texttt{simulate\_batch} 的核心是：
\begin{lstlisting}[language=Python]
exact = S0 * np.exp((MU - 0.5 * SIGMA**2) * T + SIGMA * w)
dw2 = dw * dw
base = 1 + MU * h + SIGMA * dw
em = S0 * np.prod(base, axis=1)
mil = S0 * np.prod(
    base + 0.5 * SIGMA**2 * (dw2 - h), axis=1)
endpoints = np.column_stack((em, mil))
\end{lstlisting}
\texttt{dw*dw} 是每个位置分别平方，\texttt{**2} 也是平方；它们不是矩阵乘法。\texttt{base} 的形状为 $(B,N)$，沿时间轴调用 \texttt{prod}，把每条路径的 $N$ 个乘子相乘，一次得到所有终值。数学上就是
\[
S_N=S_0F_0F_1\cdots F_{N-1}
\]
的简写，没有改变递推格式。若要画整条路径，需要 \texttt{cumprod} 保存逐步累积乘积；终值误差实验只需 \texttt{prod}。精确路径图类似地用 \texttt{cumsum} 得到各时刻的布朗运动。

\texttt{column\_stack} 把两个长度为 $B$ 的数组并成两列，所以 \texttt{endpoints.shape == (B,2)}，第 0 列为 EM，第 1 列为 Milstein。精确终值 \texttt{exact} 的形状却是 $(B,)$。作差时要写
\begin{lstlisting}[language=Python]
error = endpoints - exact[:, None]
\end{lstlisting}
\texttt{:} 表示取全部路径，\texttt{None} 插入一个新轴，令形状从 $(B,)$ 变为 $(B,1)$。NumPy 的广播机制把这一列用于两列数值解，得到 $(B,2)$ 的配对误差。它没有给每种方法重新生成一条精确路径。例如两条路径有
\[
\texttt{endpoints}=\begin{pmatrix}99&100\\110&112\end{pmatrix},
\quad \texttt{exact[:,None]}=\begin{pmatrix}101\\111\end{pmatrix},
\quad \texttt{error}=\begin{pmatrix}-2&-1\\-1&1\end{pmatrix}.
\]
这里的减法在“相同路径”内完成，这是配对估计。

\texttt{simulate\_batch} 中的 \texttt{yield} 使函数成为生成器。调用者用 \texttt{for} 逐层取出四项：步数 $N$、精确终值、两列数值终值、控制变量。程序不会一次保存每个网格的所有中间状态；所有网格共享本批已经生成的 \texttt{fine}。

\subsection{强误差、弱偏差、控制变量各占哪一列}

原程序没有把全部误差样本永久保存在内存，而是对三个指标累积总和与平方和：
\begin{lstlisting}[language=Python]
corrected = error - controls @ betas[k]
metrics = np.stack((np.abs(error), error, corrected), axis=2)
sums[k] += metrics.sum(axis=0)
squares[k] += (metrics * metrics).sum(axis=0)
\end{lstlisting}
此处 \texttt{metrics.shape == (B,2,3)}：第一个轴是路径，第二个轴是方法，第三个轴依次是绝对终值误差、带符号配对差、控制变量修正后的带符号差。跨路径求和后形状为 $(2,3)$。外层 \texttt{sums.shape == (L,2,3)}，其中 $L$ 是网格层数，\texttt{k} 指当前网格。

三列对应的总体目标依次是 $\E|S_T^h-S_T|$、$\E[S_T^h-S_T]$、以及同一个弱偏差。后两列目标相同，估计噪声不同。不能把 \texttt{np.abs(error)} 当成弱误差，也不能先将弱样本逐个取绝对值再平均。

在某一网格上，\texttt{controls.shape == (B,5)}，\texttt{betas[k].shape == (5,2)}。符号 \texttt{@} 表示矩阵乘法，故乘积是 $(B,2)$，为每条路径、每种方法分别计算五个控制变量的加权和。控制变量的严格零均值及方差缩减原理在理论部分已说明；对应代码的关键是：\emph{系数只在独立预试验中拟合，然后在正式样本中固定使用。}

\texttt{fit\_controls} 分批记录 $X$ 与 $Y$ 的总和、$X^TX$ 和 $X^TY$。$X$ 的五列是控制变量，$Y$ 的两列是 EM/Milstein 配对误差。总预试验样本数记为 $M_0$，最终进行中心化最小二乘：
\begin{lstlisting}[language=Python]
covx = sxx[k] - np.outer(sx[k], sx[k]) / count
covxy = sxy[k] - np.outer(sx[k], sy[k]) / count
betas[k] = np.linalg.solve(covx, covxy)
\end{lstlisting}
\texttt{outer} 是外积；例如两个长度分别为 5 和 2 的向量，外积为 $5\times2$ 矩阵。这里 \texttt{covx} 与 \texttt{covxy} 是尚未除以 $M_0-1$ 的中心化交叉乘积；同一个分母会在最小二乘方程两边抵消。\texttt{solve(A,B)} 求解 $AZ=B$，不必显式计算矩阵逆。预试验点数过少会使拟合不稳定，不能随意用几条路径代替原来的 20,000 条预试验路径。

原代码分别使用 \texttt{default\_rng(SEED+1)} 和 \texttt{default\_rng(SEED)} 建立预试验与正式实验的随机流。拟合阶段的中心化用于估计协方差；正式校正时使用理论已知为零均值的控制变量本身，不是用正式样本估计一个均值后再从每个控制变量中扣除。

\subsection{批量累计：平均数和标准误怎样得到}

如果 $V_1,\ldots,V_M$ 是不同路径上的某个指标，只需保存
\[
M,\qquad A=\sum_{j=1}^M V_j,\qquad B_2=\sum_{j=1}^M V_j^2,
\]
就能计算
\begin{equation}
\overline V=\frac{A}{M},\qquad
s^2=\frac{B_2-A^2/M}{M-1},\qquad
\operatorname{SE}(\overline V)=\frac{s}{\sqrt M}.
\label{code:streaming}
\end{equation}
\texttt{gbm.py} 的 \texttt{mean\_se(total,squares,count)} 把这三个公式一次应用到所有网格、方法、指标。\texttt{np.maximum(...,0)} 避免理论上非负的方差因浮点舍入产生很小的负数；它不能修复错误的数据。

非仿射程序把同一逻辑封装进 \texttt{Moments} 对象，其接口是：
\begin{lstlisting}[language=Python]
strong[n, r].add(np.abs(terminal[n] - terminal[r]))
weak[n, r].add(payoff[n] - payoff[r])
result = strong[n, r].summary()
\end{lstlisting}
\texttt{add(values)} 接收一批独立路径的观测值，更新样本数、总和与平方和；\texttt{summary()} 返回均值、标准误、95\% 区间上下端点等字段。上式字典的键 \texttt{(n,r)} 表示“粗网格 $n$ 步，对照参考网格 $r$ 步”。元组键不是乘法。

这样分批的统计结果与把所有样本放在一起平均等价，只有浮点加法顺序带来的微小差别。分批还能限制布朗增量数组的内存占用。不过原项目为绘图和事后检查，另外保留了部分或全部终值样本；它并不是完全不存样本。若数据均值极大、方差极小，式 \eqref{code:streaming} 的两大数相减会损失精度，可改用 Welford 算法；现有程序采用的是总和与平方和形式。

跨网格数据相关，因为它们共用布朗路径；同一网格上不同路径仍是 Monte Carlo 平均所用的独立观测。因此逐点标准误有明确意义，但各层 95\% 区间不是“所有层同时覆盖真值的 95\% 区间”。

\subsection{非仿射模型：相关噪声和同时更新}

\texttt{nonaffine.py} 首先生成两个独立标准正态数组，然后在线性组合中加入相关性：
\begin{lstlisting}[language=Python]
dw = rng.standard_normal((finest, m, 2))
dw[:, :, 1] = (
    PARAMETERS["rho"] * dw[:, :, 0]
    + math.sqrt(1 - PARAMETERS["rho"]**2) * dw[:, :, 1])
dw *= math.sqrt(hfinest)
\end{lstlisting}
赋值右侧计算时，\texttt{dw[:,:,1]} 还是原先独立的第二个正态噪声。令 $Z_1,Z_2$ 独立标准正态，则最终的两个增量为
\[
\Delta W_1=\sqrt h Z_1,\qquad
\Delta W_2=\sqrt h\bigl(\rho Z_1+\sqrt{1-\rho^2}Z_2\bigr).
\]
直接用概率论可核对：两者方差均为 $h$，协方差为 $\rho h$，相关系数为 $\rho$。这里的第三个轴长度为 2，记录两种噪声；“路径条数很多”不等于“每条路径有很多独立噪声分量”。

\texttt{em\_terminal} 将最细网格增量按时间分组：
\begin{lstlisting}[language=Python]
finest, n_paths, _ = dw.shape
ratio = finest // n_steps
increments = dw if ratio == 1 else dw.reshape(
    n_steps, ratio, n_paths, 2).sum(axis=1)
\end{lstlisting}
如果本层已经是最细网格，直接复用原数组；否则将相邻细增量相加。原实验的所有步数都整除 8192；教学脚本另外显式检查整除关系，避免误用不相容的网格。

在计算内核里，\texttt{x} 保存 $X=\log S$，\texttt{y} 保存 $Y$。对每个时刻、每条路径，原程序用旧状态计算两个新状态：
\begin{lstlisting}[language=Python]
old_y = y[j]
g = g_min + (g_max - g_min) * sigmoid
x[j] += (mu - 0.5*g*g)*h + g*increments[step, j, 0]
y[j] += (kappa*(theta-old_y)*h
         + xi*math.sqrt(1.0+old_y*old_y)*increments[step, j, 1])
\end{lstlisting}
这里 \texttt{sigmoid} 也是在 \texttt{old\_y} 上计算的。\texttt{+=} 表示将右侧增量加到左侧旧值。代码先写 $X$ 更新、再写 $Y$ 更新，并不意味着第二个公式使用了新 $Y$；两个公式的系数都在 $Y_n$ 上取值。这正是“同时更新”的含义。若先改写 $Y$、再用新 $Y$ 计算波动率，就已经换成了另一种离散方法。

教学脚本对全部路径一次进行向量运算，写 \texttt{old\_y = y.copy()} 后再更新。这里需要 \texttt{copy()} 的原因是：\texttt{old\_y = y} 只让两个变量名指向同一个数组；后续原地修改 \texttt{y} 会同时改变它。原 Numba 内核里的 \texttt{old\_y = y[j]} 是标量读取，不存在这个数组别名问题。

逻辑函数 $\ell(y)=1/(1+e^{-y})$ 的稳定计算也值得读懂：
\begin{lstlisting}[language=Python]
if old_y >= 0:
    sigmoid = 1.0 / (1.0 + math.exp(-old_y))
else:
    exp_y = math.exp(old_y)
    sigmoid = exp_y / (1.0 + exp_y)
\end{lstlisting}
$y\geq0$ 时 $-y\leq0$，$y<0$ 时直接算 $e^y$，所以两分支都不会计算极大的正指数。两式数学上相等，但这样避免 $y$ 很负时 $e^{-y}$ 溢出。教学脚本用布尔掩码实现同样的两分支；不直接用可能同时计算两边的 \texttt{np.where}。最终才用 \texttt{np.exp(x)} 还原价格，并检查其有限性与正性。理论上的 $e^x>0$ 不会自动阻止计算机上的上溢或下溢，因此原项目也记录中间 $X,Y$ 的范围和状态诊断。

\subsection{为什么原非仿射代码使用 Numba}

原非仿射实验包含 100,000 条路径，以及 $N=8,16,32,64,128,256,2048,4096,8192$ 共九层网格。总路径步更新数为
\[
100000\,(8+16+32+64+128+256+2048+4096+8192)
=1\,484\,000\,000.
\]
纯 Python 的双重逐项循环开销较大。\texttt{@njit(cache=False)} 让 Numba 在首次调用时把内核编译为较快的机器代码；公式没有因编译而改变。\texttt{run} 先用一个很小的零增量输入预热编译，并单独记录编译耗时，再开始统计主实验时间。因此不能只比较两个脚本日志里的秒数，就宣称某种数学方法更快：样本数、网格、绘图和预试验范围本来不同。

教学脚本去掉 Numba，时间步仍循环，但在每一步使用 NumPy 同时更新整批路径。它的较小网格足以学习算法，也能在本机当前 Python/NumPy 环境中直接运行。它不会冒充原来十万路径、最细 8192 步的大规模实验。

\subsection{结果文件、日志与报告数字怎样连接}

\texttt{CSV} 是带列名的表格，例如 \texttt{gbm\_convergence.csv} 每行对应一种方法和一个网格，并分别保存强误差、原始弱偏差、控制变量弱偏差及其标准误。\texttt{JSON} 适合嵌套结构：参数、随机种子、软件版本、拟合窗口、诊断和结果可以放进同一份摘要。\texttt{NPZ} 是 NumPy 的压缩数组容器，保存真实终值样本，便于事后从样本重新计算统计。

例如原非仿射结果存为 \texttt{nonaffine\_terminals.npz}，其中键 \texttt{S\_64} 表示 64 步网格的全部终值，\texttt{S\_8192} 表示最细参考的全部终值。下面是一段\emph{只读}练习，须在项目根目录运行：
\begin{lstlisting}[language=Python]
from pathlib import Path
import numpy as np

with np.load(Path("results") / "nonaffine_terminals.npz") as data:
    coarse = data["S_64"]
    reference = data["S_8192"]
    strong = np.mean(np.abs(coarse - reference))
    difference = (np.maximum(coarse - 100.0, 0.0)
                  - np.maximum(reference - 100.0, 0.0))
    weak = difference.mean()
    se = difference.std(ddof=1) / np.sqrt(difference.size)
print(strong, weak, se)
\end{lstlisting}
\texttt{maximum(s-100,0)} 是逐元素计算看涨收益。\texttt{ddof=1} 表示样本方差用分母 $M-1$。标准误必须从\emph{配对差} \texttt{difference} 计算，不能忽略粗细样本的协方差。

\texttt{verify.py} 就包括这种保存结果的交叉检查，还用预先指定的正弦、余弦小数组比较“向量化写法”和“逐条递推写法”。这些小数组不是布朗运动样本；它们的用途是检验代数实现，不是检验随机模型的分布。它还检查常波动率情况下 log-EM 是否退化为精确 GBM。通过后写入 \texttt{results/verification.json}。该脚本会导入 \texttt{nonaffine.py}，因此也需要 Numba；它不会重新模拟全部路径，但不是完全不写文件的检查。

\texttt{build\_report\_data.py} 不再抽样，而是读 JSON/CSV，生成 \texttt{report\_values.tex} 与各表格的 \texttt{.tex}，并更新报告图。这让报告中“EM 强阶”等数字直接来自保存结果。用 \texttt{np.polyfit} 对 $\log h$ 与 $\log e$ 作一次直线拟合，返回斜率与截距；斜率对应经验幂指数。程序对非仿射弱差先检查选定窗口中的区间是否排除零；若有区间含零，\texttt{rate\_fit} 返回 \texttt{slope=None}。即使均排除零，原程序仍将有限参考下的斜率标为描述性结果。

GBM 的 \texttt{log} 函数一边打印、一边把内容加入列表，最后写入 \texttt{gbm\_run.log}。非仿射程序打印进度和摘要，并把时间、种子与环境写进 JSON；\texttt{run\_all.py} 本身没有把所有终端输出自动收集成总日志。需要记录个人演示输出时，可明确重定向到讲义目录中的新文件。

\subsection{从小练习开始，再决定是否完整复现实验}

本讲义新增 \texttt{tutorial\_zh/teaching\_demo.py}。在项目根目录打开终端，依次运行：
\begin{lstlisting}[language=bash]
python3 tutorial_zh/teaching_demo.py --self-check
python3 tutorial_zh/teaching_demo.py
python3 tutorial_zh/teaching_demo.py --paths 8000 --seed 2026090804
\end{lstlisting}
第一条只执行确定性检查，包括 GBM 递推与精确公式、粗化后总增量一致、非仿射旧状态同时更新、常波动率退化、极端输入下逻辑函数稳定性，以及分批统计的标准误。第二条默认使用 2000 条路径、批大小 256；GBM 的层数为 $8,16,\ldots,512$，非仿射粗网格为 $16,32,64,128$，参考网格为 $256,512,1024$。第三条增加路径数，适合观察区间通常如何缩窄。脚本没有控制变量，也不自动拟合或认证收敛阶。

本次在本机 Python 3.9.6、NumPy 2.0.2 上，确定性检查通过，默认演示运行通过。例如 GBM 的 $N=512$ 时，EM 强误差约为 $0.236916$，Milstein 约为 $0.003028$；EM 的原始弱差约为 $-0.001314$，其 95\% 区间约为 $[-0.014883,0.012255]$，而解析均值偏差约为 $-0.000257$。这个小练习同时展示了“强误差明显可见”和“弱偏差可能被抽样噪声淹没”。这些是教学脚本自己的结果，不替换原报告的统计数字。

非仿射演示的参考只到 1024 步，默认样本中 $512\to1024$ 的平均绝对终值差约为 $0.116254$。与某个粗网格差相比，它可能并不小；这正好提醒我们检查参考误差，不能把最细曲线直接命名为精确解。更换种子后个别区间或误差大小会变化，不应要求每次结果都整齐地单调下降。

原项目完整运行使用 README 中记录的依赖环境。其 \texttt{requirements.txt} 固定 NumPy、Matplotlib 和 Numba 版本；已有实验元数据来自另一套环境，因此本机能运行 NumPy 教学版，并不表示已复现原环境。准备好兼容依赖、并在项目的副本中操作时，完整实验命令为：
\begin{lstlisting}[language=bash]
python3 run_all.py
\end{lstlisting}
或者分别执行：
\begin{lstlisting}[language=bash]
python3 gbm.py --paths 160000 --pilot-paths 20000 --batch-size 2000
python3 nonaffine.py --paths 100000 --batch-size 512 --seed 2026090804
python3 verify.py
python3 build_report_data.py
\end{lstlisting}
原 GBM 的种子写在常量 \texttt{SEED} 中，没有 \texttt{--seed} 命令行选项。非仿射与教学脚本有该选项。原报告源是英文，可按其 README 用 \texttt{pdflatex} 编译；本中文讲义采用支持中文的引擎，编译方法见随讲义生成的说明。

最后要区分“同分布复现实验”和“尽量得到同一串输出数字”。固定种子决定伪随机流，但还应保留生成器类型、软件版本、路径数、网格、批大小及抽样调用顺序。非仿射数组是 $(N_f,B,2)$；改变 $B$ 后，同一串随机数会分配到不同的路径和时间位置，所以即使种子相同，逐路径终值仍会改变。GBM 与非仿射在原项目中使用相同种子，也不意味着二者共享可逐路径比较的噪声；它们是两个不同实验。教学脚本有意分别使用传入的种子及其加一值，并将它们打印出来。

建议实际练习时先读清一批的数组形状，再手算一条两步路径，接着核对默认输出中的强误差与带符号弱差，最后查看保存终值能否还原原报告统计。这样每一层代码都有可解释、可核对的数学对象。
